% Einheit 1 von 5 – Winkel messen und zeichnen (bank/winkel-dreiecke/e1.jsonl, zusammenbau v0.1)
%% TODO Zweigzeile Teil 1: was hier gelernt wird (Satz in Schülersprache aus dem Einheitstitel) – Einheit 1
\einheitenkopf[e1]{Einheit 1 von 5 · Winkel messen und zeichnen}
\zweigzeile{ab Kl. 5 · P10}
\setcounter{aufgabe}{7}

%% TODO Titel: Ich-kann-Satz fehlt in der Bank (Platzhalter: Kettenname) – Nr. 8
%% TODO Anweisung: ein Satz über der ersten Teilaufgabe fehlt in der Bank – Nr. 8
\begin{aufgabe}{Welche Skala?}
\begin{teile}
\teil Nicht ablesen: Gilt beim Messen von α am Geodreieck die kleinere oder die größere der beiden Zahlen? \\ \kreuz{die kleinere Zahl} \\ \kreuz{die größere Zahl}
\end{teile}
\end{aufgabe}

%% TODO Titel: Ich-kann-Satz fehlt in der Bank (Platzhalter: Kettenname) – Nr. 9
%% TODO Anweisung: ein Satz über der ersten Teilaufgabe fehlt in der Bank – Nr. 9
\begin{aufgabe}{Winkel messen}
\begin{teile}
\teil Welche Winkelart hat α? Nicht messen. \\ \kreuz{spitz} \\ \kreuz{recht} \\ \kreuz{stumpf} \\ \kreuz{gestreckt} \\ \kreuz{überstumpf}
\end{teile}
\swz{Miss den Winkel α mit dem Geodreieck.}{\winkel[5]{37}{\alpha}}
\swz{Miss den Winkel α mit dem Geodreieck.}{\winkel[5]{54}{\alpha}}
\swz{Miss den Winkel α mit dem Geodreieck.}{\winkel[5]{18}{\alpha}}
\swz{Miss den Winkel α mit dem Geodreieck.}{\winkel[5]{83}{\alpha}}
\swz{Miss den Winkel α mit dem Geodreieck.}{\winkel[5]{66}{\alpha}}
\swfrage{Was bleibt gleich, was ändert sich?}
\swz{Miss den stumpfen Winkel α mit dem Geodreieck.}{\winkel[5]{118}{\alpha}}
\swa{Zeichne an den Strahl mit dem Scheitel S einen Winkel von $55^\circ$.}{\winkelstrahl}
\swz{Bestimme den überstumpfen Winkel α. Miss dazu den Rest bis zum Vollwinkel.}{\winkel[4]{215}{\alpha}}
\swz[3]{Schreibe den Winkel α im Dreieck ABC mit drei Buchstaben ($\angle$…). Welcher Punkt ist der Scheitel?}{\dreieck{(0,0)}{(5,0)}{(1.5,3)}{a}{b}{c}{\alpha}{\beta}{\gamma}}
%% TODO Darstellung für schwach fehlt in der Bank (winkel-dreiecke-e1-k2-s6-v1; Abschnitt „Für schwache Schüler“)
\swz{Die Winkel $34^\circ$ und $47^\circ$ liegen am selben Scheitel nebeneinander. Wie groß ist der Winkel, den beide zusammen bilden?}{}
%% TODO Darstellung für schwach fehlt in der Bank (winkel-dreiecke-e1-k2-s7-v1; Abschnitt „Für schwache Schüler“)
\swz{Ein Winkel von $96^\circ$ ist durch einen Strahl in zwei Teilwinkel geteilt. Der eine misst $38^\circ$. Wie groß ist der andere?}{}
%% TODO Darstellung für schwach fehlt in der Bank (winkel-dreiecke-e1-k2-s8-v1; Abschnitt „Für schwache Schüler“)
\swz{Ein Wanderweg ist 1,2 km lang. Nach 450 m macht Ella eine Pause. Wie viele Meter hat sie noch vor sich?}{}
\swz[3]{Vom Punkt A aus sieht Jana die Spitze T eines Turms unter einem Winkel von $41^\circ$ zur waagerechten Grundlinie, die Spitze F der Fahne auf dem Turm unter $44^\circ$. Wie groß ist der Winkel γ zwischen den Sichtlinien AT und AF? \hfill (P10 2018 OS)}{}
\end{aufgabe}

%% TODO Titel: Ich-kann-Satz fehlt in der Bank (Platzhalter: Kettenname) – Nr. 10
%% TODO Anweisung: ein Satz über der ersten Teilaufgabe fehlt in der Bank – Nr. 10
\begin{aufgabe}{Winkel schätzen}
\begin{teile}
\teil Schätze, ohne zu messen: Wie groß ist α im Vergleich zu einem rechten Winkel? \\ \kreuz{ein Viertel} \\ \kreuz{die Hälfte} \\ \kreuz{drei Viertel}
\end{teile}
\end{aufgabe}

%% TODO Titel: Ich-kann-Satz fehlt in der Bank (Platzhalter: Kettenname) – Nr. 11
%% TODO Anweisung: ein Satz über der ersten Teilaufgabe fehlt in der Bank – Nr. 11
\begin{aufgabe}{Winkel messen – Fehler finden, Begründen, Anwendung}
\swz[3]{Tim hat den Winkel α gemessen und notiert: $\alpha = 144^\circ$. Finde den Fehler und gib den richtigen Wert an.}{\winkel[5]{36}{\alpha}}
\swfrage{Warum heißt ein Winkel von $91^\circ$ stumpf, ein Winkel von $89^\circ$ aber spitz? Begründe.}
\swz[3]{Eine Parkhausschranke dreht sich beim Öffnen aus der waagerechten Lage, bis sie senkrecht steht. Nach zwei Sekunden hat sie sich um $54^\circ$ gedreht. Um wie viel Grad muss sie sich noch drehen?}{}
\end{aufgabe}

\uebersichtskasten{Winkel: zwei Schenkel und ein Scheitel; gemessen wird die Drehung vom einen Schenkel zum anderen, in Grad. \\ spitz: unter 90° \quad recht: 90° \quad stumpf: zwischen 90° und 180° \quad gestreckt: 180° \quad überstumpf: über 180° \quad Vollwinkel: 360° \\ Messen: Nullpunkt des Geodreiecks auf den Scheitel, Nulllinie auf einen Schenkel, am anderen Schenkel ablesen – beim spitzen Winkel die kleinere Zahl, beim stumpfen die größere. \\ Winkel über 180°: den Rest messen und von 360° abziehen: 360° − 140° = 220° \\ Teilwinkel: nebeneinander → addieren; einer im anderen → abziehen. \quad α = 25° + 40° = 65° \quad β = 125° − 45° = 80°}
